\documentclass[
    language=english,
    doctype=manual,
    institution=none,
    titlestyle=book,
    oneside,
]{../../omnilatex}

\RequirePackage{config/document-settings}

\addbibresource{../../bib/bibliography.bib}

\begin{document}

\maketitle

\tableofcontents

%=============================================================================
\chapter{Overview}
\label{sec:overview}
%=============================================================================

This runbook provides operational procedures for the DataPipe production
environment. It covers monitoring, incident response, routine
maintenance, and escalation paths. All on-call engineers must be familiar
with the contents of this manual before their first rotation.

\section{Environment Topology}

The production environment spans three availability zones in the
us-east-1 region:

\begin{figure}[htbp]
    \centering
    \begin{tikzpicture}[
        node distance=1.2cm and 2cm,
        box/.style={draw, rounded corners, minimum width=2.6cm,
                     minimum height=0.9cm, text centered, font=\small,
                     fill=#1!20, draw=#1!60!black},
        az/.style={draw, dashed, rounded corners, inner sep=8pt,
                    minimum width=3.5cm, minimum height=3.5cm,
                    label={[font=\footnotesize]above:#1}},
        arr/.style={->, >=stealth, thick},
    ]
        \begin{scope}[local bounding box=az-a]
            \node[box=blue]   (k1) {Kafka-1};
            \node[box=blue]   (d1) [below=of k1] {DataPipe-1};
            \node[box=blue]   (p1) [below=of d1] {PostgreSQL-1};
        \end{scope}
        \node[az, fit=az-a, draw=blue!40, label={[font=\footnotesize]above:AZ-a}] {};

        \begin{scope}[xshift=5.5cm, local bounding box=az-b]
            \node[box=green]  (k2) {Kafka-2};
            \node[box=green]  (d2) [below=of k2] {DataPipe-2};
            \node[box=green]  (p2) [below=of d2] {PostgreSQL-2};
        \end{scope}
        \node[az, fit=az-b, draw=green!40, xshift=5.5cm,
              label={[font=\footnotesize]above:AZ-b}] {};

        \begin{scope}[xshift=11cm, local bounding box=az-c]
            \node[box=orange] (k3) {Kafka-3};
            \node[box=orange] (d3) [below=of k3] {DataPipe-3};
            \node[box=orange] (p3) [below=of d3] {PostgreSQL-3};
        \end{scope}
        \node[az, fit=az-c, draw=orange!40, xshift=11cm,
              label={[font=\footnotesize]above:AZ-c}] {};

        \draw[arr] (k1) -- (k2);
        \draw[arr] (k2) -- (k3);
        \draw[arr] (k3) -- (k1);
    \end{tikzpicture}
    \caption{Production environment topology across three availability zones.}
    \label{fig:topology}
\end{figure}

%=============================================================================
\chapter{Monitoring Procedures}
\label{sec:monitoring}
%=============================================================================

\section{Dashboards}

All dashboards are available in Grafana at \texttt{https://grafana.internal}.

\begin{table}[htbp]
    \centering
    \caption{Grafana dashboards and their purpose.}
    \label{tab:dashboards}
    \begin{tabular}{@{}lp{8cm}@{}}
        \toprule
        Dashboard & Purpose \\
        \midrule
        DataPipe Overview & High-level pipeline health, throughput, and error rates. \\
        Kafka Consumer Lag & Per-consumer-group lag metrics and alerts. \\
        PostgreSQL Performance & Query latency, connection count, replication lag. \\
        Infrastructure & CPU, memory, disk, and network across all nodes. \\
        \bottomrule
    \end{tabular}
\end{table}

\section{Health Checks}

Run the following commands during routine monitoring:

\begin{lstlisting}[language=bash, caption={Health check commands.}]
# Cluster-wide health check
curl -s https://datapipe.internal/api/v1/cluster/health | jq .

# Individual node check
ssh datapipe-node-01 "dp cluster status"

# Kafka consumer lag check
kafka-consumer-groups.sh --bootstrap-server kafka:9092 \\
    --describe --group datapipe-etl
\end{lstlisting}

\section{Key Metrics}

Monitor the following metrics continuously:

\begin{itemize}
    \item \textbf{Pipeline throughput} (records/second) --- should remain
        above 50\,000 rps in production.
    \item \textbf{Consumer lag} --- alert if exceeds 10\,000 records for
        more than 5 minutes.
    \item \textbf{Error rate} --- alert if exceeds 0.1\% of total records.
    \item \textbf{End-to-end latency} --- alert if P99 exceeds 2 seconds.
    \item \textbf{Disk usage} --- alert at 80\% capacity.
    \item \textbf{Replication lag} --- alert if exceeds 100\,ms.
\end{itemize}

%=============================================================================
\chapter{Incident Response}
\label{sec:incident-response}
%=============================================================================

\section{Severity Levels}

\begin{table}[htbp]
    \centering
    \caption{Incident severity classification.}
    \label{tab:severity}
    \begin{tabular}{@{}lp{5cm}p{3cm}@{}}
        \toprule
        Level & Description & Response Time \\
        \midrule
        SEV-1 & Complete data loss or pipeline failure & 15 minutes \\
        SEV-2 & Partial data loss or degraded throughput & 30 minutes \\
        SEV-3 & Non-critical error or cosmetic issue & 4 hours \\
        SEV-4 & Informational or low-impact issue & 24 hours \\
        \bottomrule
    \end{tabular}
\end{table}

\section{Escalation Matrix}

\begin{table}[htbp]
    \centering
    \caption{Escalation contacts by severity and time.}
    \label{tab:escalation}
    \begin{tabular}{@{}llll@{}}
        \toprule
        Level & First Responder & After 30 min & After 1 hour \\
        \midrule
        SEV-1 & On-call engineer & SRE lead & VP Engineering \\
        SEV-2 & On-call engineer & SRE lead & SRE lead \\
        SEV-3 & On-call engineer & Team lead & SRE lead \\
        SEV-4 & On-call engineer & Team lead & --- \\
        \bottomrule
    \end{tabular}
\end{table}

\section{Incident Response Playbook}

\subsection{Pipeline Failure (SEV-1)}

\begin{enumerate}
    \item Acknowledge the alert in PagerDuty.
    \item Assess the scope: is the failure affecting one pipeline or all
        pipelines?
    \item If a single pipeline:
    \begin{enumerate}
        \item Check recent deployments or configuration changes.
        \item Review pipeline logs: \texttt{dp logs --pipeline <name> --level error}.
        \item Attempt a graceful restart: \texttt{dp pipeline restart <name>}.
    \end{enumerate}
    \item If all pipelines:
    \begin{enumerate}
        \item Check cluster health: \texttt{dp cluster health}.
        \item Verify Kafka connectivity.
        \item Check for infrastructure-level issues (disk, network).
    \end{enumerate}
    \item If restart does not resolve, escalate to SRE lead.
    \item Document the incident in the incident tracker.
\end{enumerate}

\subsection{Kafka Broker Down (SEV-2)}

\begin{enumerate}
    \item Verify which broker is affected via the Kafka dashboard.
    \item Check if replication factor is sufficient (minimum 3).
    \item Confirm remaining brokers have sufficient capacity.
    \item Monitor consumer lag during recovery.
    \item If lag grows beyond threshold, scale up consumers temporarily.
    \item After broker recovery, verify partition rebalancing completes.
\end{enumerate}

\subsection{Database Replication Lag (SEV-2)}

\begin{enumerate}
    \item Check PostgreSQL replication status:
    \begin{lstlisting}[language=bash]
psql -h primary -c "SELECT client_addr, sent_lsn,
    replay_lag FROM pg_stat_replication;"
    \end{lstlisting}
    \item If lag exceeds 1 minute, check for long-running queries on
        replicas.
    \item Kill any blocking queries and monitor recovery.
    \item If lag persists, consider rebuilding the replica from a backup.
\end{enumerate}

%=============================================================================
\chapter{Routine Maintenance}
\label{sec:maintenance}
%=============================================================================

\section{Daily Tasks}

\begin{enumerate}
    \item Review overnight monitoring dashboards for anomalies.
    \item Check consumer lag across all pipelines.
    \item Verify backup completion and integrity.
    \item Review error logs for recurring issues.
\end{enumerate}

\section{Weekly Tasks}

\begin{enumerate}
    \item Review capacity planning metrics and projections.
    \item Audit access logs and API usage.
    \item Test disaster recovery procedures.
    \item Review and update runbook entries.
\end{enumerate}

\section{Monthly Tasks}

\begin{enumerate}
    \item Apply security patches and updates.
    \item Review and rotate API keys and credentials.
    \item Capacity planning review with engineering leadership.
    \item Post-incident review for all SEV-1 and SEV-2 incidents.
\end{enumerate}

\section{Backup and Recovery}

\begin{lstlisting}[language=bash, caption={Backup commands.}]
# Full cluster backup
dp backup create --full --output s3://backups/datapipe/

# Configuration backup
dp backup config --output s3://backups/datapipe/config/

# Verify backup integrity
dp backup verify s3://backups/datapipe/latest/
\end{lstlisting}

%=============================================================================
\chapter{Contact Information}
\label{sec:contacts}
%=============================================================================

\begin{table}[htbp]
    \centering
    \caption{Emergency contacts.}
    \label{tab:contacts}
    \begin{tabular}{@{}lll@{}}
        \toprule
        Role & Name & Contact \\
        \midrule
        SRE Lead & Dr. Mira Patel & +1-555-0101 \\
        On-Call Engineer & Rotation & PagerDuty \\
        Database Admin & Jordan Lee & +1-555-0102 \\
        Security Lead & Priya Nair & +1-555-0103 \\
        \bottomrule
    \end{tabular}
\end{table}

\printbibliography[heading=bibintoc, title={References}]

\end{document}
